\documentclass[9pt]{extarticle}
\usepackage[margin=1.5in]{geometry}

% math packages
\usepackage{amsmath}
\usepackage{amsfonts}
\usepackage{mathtools}

% graphics and figures
\usepackage{graphicx}
\usepackage{tikz}
\usepackage{multirow}

% code listings
\usepackage{listings}
\lstdefinestyle{default}{basicstyle=\ttfamily\small, commentstyle=, columns=flexible, keepspaces=true, tabsize=2, breaklines=true, breakatwhitespace=true, numbers=left, numberstyle=\tiny, xleftmargin=5mm, captionpos=t}
\lstset{style=default}

% captions
\usepackage{caption}
\captionsetup{font={small, it}}

% quotes
\usepackage[english]{babel}
\usepackage[autostyle, english = american]{csquotes}
\MakeOuterQuote{"}

% hyperlinks
\usepackage{hyperref}
\hypersetup{hidelinks}

% lists and enumeration
\usepackage{enumitem}

% text spacing
\usepackage{parskip}
\usepackage{setspace}
\setlength{\parskip}{10pt}
\setstretch{1.05}

% align env. spacing
\setlength{\jot}{10pt}

% metadata
\title{Assignment}
\author{Aarush Kumbhakern}
\date{\today}

\begin{document}
\maketitle

\subsection{Experiments on XML Inputs}

To test minimisation on highly structured data, we revisit the XML predicates from the Zhang et al. probabilistic delta debugging study (CITE). Each predicate pairs an XQuery program with a bug-inducing XML document whose nested tags, cross-referenced attributes, and whitespace-sensitive fragments are render byte-oriented reducers susceptible to breakage. This setting offers a stringent check on how quickly a minimiser can prune redundant markup without destroying the semantics that trigger the failure.

The benchmark records three main metrics per run: wall-clock time, the number of oracle invocations, and the reduction ratio $(|x| - |x'|)/|x|$. Across all XML predicates tested, ZipMin consistently finishes $\approx 9\%$ sooner and issues $\approx 22\%$ fewer oracle calls than ddmin, indicating that the sweep phase quickly filters out uninteresting spans. The trade-off is a marginally lower reduction ratio, with ZipMin occasionally leaving a small sentinel fragment. In practice the difference remains within rounding error.

Benchmarking hardware: (TODO). A repository containing the self-sufficient benchmarking setup is available for reproduction\footnote{\url{https://github.com/Positron11/dd-zipmin}}.
 
\begin{table}[ht]
    \centering
    \scriptsize
    
    \caption{Benchmark results on XML predicates (predicate IDs include variant suffixes).}
    
    \begin{tabular}{l r r r r r r r}
        \hline
        \multirow{2}{*}{Predicate} & \multirow{2}{*}{Input bytes} & \multicolumn{2}{c}{Oracle Calls} & \multicolumn{2}{c}{Wall Time (h)} & \multicolumn{2}{c}{Reduction Ratio} \\
        &  & ddmin & ZipMin & ddmin & ZipMin & ddmin & ZipMin \\
        \hline
        xml-1e9bc83-1:1 & 5444 & 8260 & 6257 & 2.52 & 2.15 & 0.287 & 0.294 \\
        xml-1e9bc83-2:1 & 6965 & 10227 & 8457 & 3.20 & 2.95 & 0.297 & 0.300 \\
        xml-1e9bc83-3:1 & 7939 & 11065 & 9175 & 3.33 & 3.10 & 0.301 & 0.303 \\
        xml-1e9bc83-4:1 & 9133 & 12711 & 9543 & 3.56 & 3.27 & 0.234 & 0.242 \\
        xml-1e9bc83-5:1 & 10204 & 15132 & 11104 & 3.67 & 3.38 & 0.257 & 0.264 \\
        \hline
    \end{tabular}
    
    \label{tab:xml}
\end{table}

Overall, the XML study shows that mildly structure-aware heuristics can deliver a clear runtime and oracle-count advantage even without a full hierarchical model. The gap widens as the XML fixtures grow: for the largest predicate (\verb|xml-1e9bc83-5:1|) ZipMin trims the runtime by about $17$ minutes while sparing more than $4{,}000$ oracle invocations. The paired rows in Table~\ref{tab:xml} also underline that ZipMin never regresses relative to ddmin on these workloads.

\end{document}
